\documentclass[14pt,a4paper]{extarticle}

\usepackage[T2A]{fontenc}
\usepackage[utf8]{inputenc}
\usepackage[russian]{babel}
\usepackage{cmap}
\usepackage{geometry}
\usepackage{array}
\usepackage{longtable}
\usepackage{xcolor}
\usepackage{ragged2e}
\usepackage{listings}
\usepackage{enumitem}
\usepackage{hyperref}

\geometry{left=30mm,right=15mm,top=20mm,bottom=20mm}
\hypersetup{colorlinks=true,linkcolor=black,urlcolor=blue}
\setlength{\parindent}{1.25cm}
\setlength{\parskip}{0.2em}
\renewcommand{\arraystretch}{1.35}
\sloppy

\newcolumntype{L}[1]{>{\RaggedRight\arraybackslash}p{#1}}
\newcommand{\code}[1]{\texttt{\detokenize{#1}}}
\newcommand{\pathcode}[1]{{\footnotesize\nolinkurl{#1}}}

\lstdefinestyle{codestyle}{
    basicstyle=\ttfamily\small,
    breaklines=true,
    frame=single,
    columns=fullflexible,
    keepspaces=true,
    showstringspaces=false,
    backgroundcolor=\color{gray!5}
}

\begin{document}

\begin{titlepage}
    \begin{center}
        САНКТ-ПЕТЕРБУРГСКИЙ НАЦИОНАЛЬНЫЙ\\
        ИССЛЕДОВАТЕЛЬСКИЙ УНИВЕРСИТЕТ ИТМО

        \vspace{1.5cm}
        Дисциплина: Бэк-энд разработка

        \vspace{1.5cm}
        \textbf{Отчет}

        \vspace{0.5cm}
        Лабораторная работа 4

        \vspace{0.5cm}
        \textbf{Развёртывание приложения на удалённом сервере}

        \vfill

        \begin{flushright}
            Выполнил:\\
            Цатинян Артём\\
            Группа БР1.2

            \vspace{0.8cm}
            Проверил:\\
            Добряков Д. И.
        \end{flushright}

        \vfill
        Санкт-Петербург\\
        2026 г.
    \end{center}
\end{titlepage}

\section{Задача}

Тема проекта: приложение для бронирования столиков в ресторанах.

В рамках лабораторной работы требовалось выполнить ручное развёртывание backend-приложения на удалённом сервере. В качестве разворачиваемого приложения использовалась микросервисная версия проекта: \code{identity-service}, \code{catalog-service}, \code{booking-service}, \code{review-service}, PostgreSQL и RabbitMQ.

\section{Ход работы}

\subsection{Подготовка сервера}

Для развёртывания был использован удалённый сервер с публичным IP-адресом \code{186.246.13.204}. На сервере установлена операционная система Debian GNU/Linux 13.

На сервер были установлены необходимые компоненты:

\begin{lstlisting}[style=codestyle]
apt-get update
apt-get install -y git docker.io docker-compose openjdk-21-jdk postgresql rabbitmq-server
systemctl enable --now docker postgresql rabbitmq-server
\end{lstlisting}

Docker и Docker Compose были установлены для поддержки контейнерного сценария, однако во время ручного деплоя сервер упёрся в лимиты загрузки образов из публичных registry. Поэтому для финального ручного развёртывания был выбран более стабильный вариант для небольшой VPS: запуск собранных Spring Boot jar-файлов через \code{systemd}. Такой способ соответствует ручному развёртыванию и не зависит от Docker Hub rate limit.

\subsection{Настройка базы данных и брокера сообщений}

В PostgreSQL были созданы пользователь и база данных, которые используются сервисами:

\begin{lstlisting}[style=codestyle]
CREATE ROLE storage LOGIN PASSWORD 'storage';
CREATE DATABASE storage_local OWNER storage;
\end{lstlisting}

Для RabbitMQ был включён management plugin и создан пользователь:

\begin{lstlisting}[style=codestyle]
rabbitmq-plugins enable rabbitmq_management
rabbitmqctl add_user storage storage
rabbitmqctl set_user_tags storage administrator
rabbitmqctl set_permissions -p / storage '.*' '.*' '.*'
\end{lstlisting}

\subsection{Сборка приложения}

На сервере была предпринята попытка собрать проект через Gradle, но VPS имеет около 2 GB RAM и не имеет swap. Kotlin compiler daemon создавал слишком высокую нагрузку, поэтому сборка была перенесена на локальную машину, а на сервер загружались уже готовые Spring Boot jar-файлы.

Локальная сборка выполнялась командой:

\begin{lstlisting}[style=codestyle]
./gradlew.bat --console=plain --offline bootJar
\end{lstlisting}

После сборки на сервер были загружены jar-файлы:

\begin{longtable}{|L{0.34\textwidth}|L{0.45\textwidth}|}
\hline
\textbf{Сервис} & \textbf{Jar-файл} \\
\hline
\code{identity-service} & \pathcode{identity-service/build/libs/identity-service.jar} \\
\hline
\code{catalog-service} & \pathcode{catalog-service/build/libs/catalog-service.jar} \\
\hline
\code{booking-service} & \pathcode{booking-service/build/libs/booking-service.jar} \\
\hline
\code{review-service} & \pathcode{review-service/build/libs/review-service.jar} \\
\hline
\end{longtable}

\subsection{Настройка systemd}

Для запуска микросервисов был добавлен скрипт \code{deploy-systemd.sh}. Он создаёт четыре unit-файла:

\begin{longtable}{|L{0.42\textwidth}|L{0.35\textwidth}|}
\hline
\textbf{Unit} & \textbf{Порт сервиса} \\
\hline
\pathcode{restaurant-booking-identity.service} & 8081 \\
\hline
\pathcode{restaurant-booking-catalog.service} & 8082 \\
\hline
\pathcode{restaurant-booking-booking.service} & 8083 \\
\hline
\pathcode{restaurant-booking-review.service} & 8084 \\
\hline
\end{longtable}

Для каждого сервиса задаются переменные окружения подключения к PostgreSQL и RabbitMQ. Пример unit-файла:

\begin{lstlisting}[style=codestyle]
[Service]
WorkingDirectory=/opt/restaurant-booking-lab2
Environment=SPRING_DATASOURCE_URL=jdbc:postgresql://localhost:5432/storage_local
Environment=SPRING_DATASOURCE_USERNAME=storage
Environment=SPRING_DATASOURCE_PASSWORD=storage
Environment=SPRING_RABBITMQ_HOST=localhost
ExecStart=/usr/bin/java -jar /opt/restaurant-booking-lab2/identity-service/build/libs/identity-service.jar
Restart=always
\end{lstlisting}

Сервисы включены в автозагрузку:

\begin{lstlisting}[style=codestyle]
systemctl enable restaurant-booking-identity
systemctl enable restaurant-booking-catalog
systemctl enable restaurant-booking-booking
systemctl enable restaurant-booking-review
\end{lstlisting}

\subsection{Результат развёртывания}

После запуска все четыре сервиса находятся в состоянии \code{active}. Проверка выполнялась командами:

\begin{lstlisting}[style=codestyle]
systemctl is-active restaurant-booking-identity restaurant-booking-catalog restaurant-booking-booking restaurant-booking-review
curl http://localhost:8081/v3/api-docs
curl http://localhost:8082/v3/api-docs
curl http://localhost:8083/v3/api-docs
curl http://localhost:8084/v3/api-docs
\end{lstlisting}

Также была проверена пользовательская ручка каталога:

\begin{lstlisting}[style=codestyle]
curl http://localhost:8082/api/v1/cuisines
\end{lstlisting}

Ответ содержал тестовые кухни:

\begin{lstlisting}[style=codestyle]
[
  {"id":1,"name":"European"},
  {"id":2,"name":"Italian"},
  {"id":3,"name":"Japanese"},
  {"id":4,"name":"Vegetarian"}
]
\end{lstlisting}

\subsection{Настройка nginx reverse proxy}

По полному условию лабораторной работы также был настроен nginx в режиме обратного прокси. На сервер уже был направлен домен \code{ezhidze.ru}, поэтому приложение было опубликовано через HTTPS на этом домене. Старый активный конфиг сайта был сохранён в резервную копию, а доменный конфиг nginx заменён на reverse proxy для микросервисов.

Корневой маршрут перенаправляет пользователя на Swagger каталога:

\begin{lstlisting}[style=codestyle]
location = / {
    return 302 /catalog/swagger-ui/index.html;
}
\end{lstlisting}

Для сервисов настроены отдельные префиксы:

\begin{lstlisting}[style=codestyle]
location /identity/ { proxy_pass http://127.0.0.1:8081/; }
location /catalog/  { proxy_pass http://127.0.0.1:8082/; }
location /booking/  { proxy_pass http://127.0.0.1:8083/; }
location /review/   { proxy_pass http://127.0.0.1:8084/; }
\end{lstlisting}

После настройки выполнялась проверка:

\begin{lstlisting}[style=codestyle]
nginx -t
systemctl reload nginx
curl https://ezhidze.ru/catalog/api/v1/cuisines
curl https://ezhidze.ru/identity/v3/api-docs
curl https://ezhidze.ru/catalog/v3/api-docs
curl https://ezhidze.ru/booking/v3/api-docs
curl https://ezhidze.ru/review/v3/api-docs
\end{lstlisting}

После развёртывания приложение доступно через nginx:

\begin{longtable}{|L{0.28\textwidth}|L{0.55\textwidth}|}
\hline
\textbf{Сервис} & \textbf{Swagger UI} \\
\hline
\code{identity-service} & \pathcode{https://ezhidze.ru/identity/swagger-ui/index.html} \\
\hline
\code{catalog-service} & \pathcode{https://ezhidze.ru/catalog/swagger-ui/index.html} \\
\hline
\code{booking-service} & \pathcode{https://ezhidze.ru/booking/swagger-ui/index.html} \\
\hline
\code{review-service} & \pathcode{https://ezhidze.ru/review/swagger-ui/index.html} \\
\hline
\end{longtable}

Прямые порты сервисов \code{8081--8084} также остаются доступными для отладки, но основной внешний доступ для защиты организован через nginx и домен \code{ezhidze.ru}.

\section{Вывод}

В рамках лабораторной работы микросервисное backend-приложение было вручную развёрнуто на удалённом сервере. На сервере установлены Java 21, PostgreSQL, RabbitMQ и nginx, настроены база данных и брокер сообщений, а четыре Spring Boot сервиса запущены как systemd-сервисы с автоперезапуском и автозагрузкой.

Результат проверен через Swagger/OpenAPI endpoints и тестовую ручку каталога. Приложение доступно извне через nginx reverse proxy по домену \code{ezhidze.ru}.

\end{document}
